%===========================================
% ESTE ARQUIVO É UMA INSTRUÇÃO/PROMPT DETALHADO PARA INTELIGÊNCIA ARTIFICIAL.
%
% OBJETIVO: Criar um template LaTeX COMPILÁVEL, TOTALMENTE MODULAR e livre de "gambiarras",
% utilizando EXCLUSIVAMENTE a classe MEMOIR para atender aos padrões ABNT e aos
% requisitos de numeração e estilo do usuário.
%
% OBSERVAÇÃO CRÍTICA: A modularização (criação de macros, \makechapterstyle, etc.)
% DEVE ser feita diretamente neste arquivo .tex, NÃO em um arquivo de classe (.cls)
% separado, para facilitar os testes e a depuração.
%===========================================
\documentclass[12pt,a4paper,oneside]{memoir}
\usepackage[brazilian]{babel}
\usepackage{fontspec}
\usepackage{tikz} % NECESSÁRIO PARA O ESTILO DE CAPÍTULO PERSONALIZADO
%===========================================

\setmainfont{CMU Serif}
\setlrmarginsandblock{3cm}{2cm}{*}
\setulmarginsandblock{3cm}{2cm}{*}
\checkandfixthelayout
\OnehalfSpacing % Espaçamento entre linhas de 1,5
\setlength{\parindent}{1.25cm} % Recuo do paragrafo
%===========================================

% ===========================================
% IMPLEMENTAÇÃO IA: ESTILO DE CAPÍTULO NÃO NUMERADO (ABNT SIMPLIFICADO)
% Título CENTRALIZADO, COM LINHA HORIZONTAL (ABNT), SEM O NÚMERO GRANDE.
% ===========================================
\makechapterstyle{estiloNaoNumerado}{
    \setlength{\beforechapskip}{0pt} % Espaço antes do título
    \renewcommand{\printchaptername}{}
    \renewcommand{\chapternamenum}{}
    \renewcommand{\afterchapternum}{\par\nobreak\vskip 1.3cm}
    \renewcommand{\printchapternum}{}
    \renewcommand{\printchaptertitle}[1]{%
        \centering
        \vspace*{0pt} % Ajuste vertical
        \rule{\linewidth}{0.3pt}\par % Linha horizontal
        \vskip 1em % Espaçamento entre a linha e o texto
        \MakeUppercase{\bfseries\huge ##1}\par
        \vskip 1em % Espaçamento após o título
        \rule{\linewidth}{0.3pt}\par % Linha horizontal
    }
}

% ===========================================
% ESTILO NUMERADO PRINCIPAL (TikZ) - SEM ALTERAÇÃO
% ===========================================
\makechapterstyle{CapitulosNumerado}{
    \setlength{\beforechapskip}{0.46cm}
    \renewcommand{\printchaptername}{}
    \renewcommand{\chapternamenum}{}
    \renewcommand{\afterchapternum}{\par\nobreak\vskip 1.3cm}
    \renewcommand{\printchapternum}{%
        \noindent
        \begin{tikzpicture}[remember picture, overlay]
            \def\numsize{61}
            \def\barwidth{0.8pt}
            \def\leftoffset{2.3cm}

            \node[inner sep=-1pt, anchor=south west, opacity=0] (num) at (\leftoffset, 0)
             {\fontsize{\numsize}{\numsize}\selectfont\rmfamily\thechapter};
            \node[inner sep=0pt, anchor=south west] at (\leftoffset, -1pt)
             {\fontsize{\numsize}{\numsize}\selectfont\rmfamily\thechapter};
            \node[inner sep=0pt, anchor=south west, font=\Large\rmfamily] (chaptext)
             at (0pt, -4pt) {\chaptername};

            \coordinate (StartPoint) at (0.5*\barwidth, 13pt);
            \draw[line width=\barwidth] (StartPoint) -- (StartPoint |- num.north) -- (num.north west);
            \coordinate (RightWall) at (\linewidth, 0);
            \draw[line width=\barwidth] (num.north east) -- (RightWall |- num.north) -- (RightWall |- num.south) -- (num.south east);
        \end{tikzpicture}%
    }
    \renewcommand{\printchaptertitle}[1]{%
        \raggedright\huge\bfseries ##1
    }
}

% ===========================================
% IMPLEMENTAÇÃO IA: COMANDO MODULAR DE CAPÍTULO CENTRALIZADO
% Usa \chapter* e adiciona a entrada no Sumário (TOC)
% ===========================================
\newcommand{\capituloCentralizado}[1]{
    \chapterstyle{estiloNaoNumerado} % Ativa o estilo simplificado
    \chapter*{#1}
    \addcontentsline{toc}{chapter}{#1} % Adiciona no Sumário
    \chapterstyle{CapitulosNumerado} % Retorna ao estilo principal
}
%===========================================
% DEFINIÇÕES AUXILIARES (SEM ALTERAÇÃO)
%===========================================
\newcommand{\folha}{
    \thispagestyle{empty}
    \vspace*{8cm}
    \begin{center}
        {\bfseries\scshape\Huge Folha de Rosto \par}
    \end{center}
    \clearpage
}
%===========================================
\newcommand{\contra}{
    \thispagestyle{empty}
    \vspace*{8cm}
    \begin{center}
        {\bfseries\scshape\Huge Contra Capa \par}
    \end{center}
    \clearpage
}
%===========================================

% ===========================================
% IMPLEMENTAÇÃO IA: CONFIGURAÇÃO DE CABEÇALHOS (RUNNING HEADS) E MARCA
% Garante que \thispagestyle{plain} seja aplicado à primeira página do capítulo.
% ===========================================
\makepagestyle{meuestilo}
\makeevenhead{meuestilo}{\footnotesize\textit{\leftmark}}{}{\footnotesize\thepage}
\makeoddhead {meuestilo}{\footnotesize\textit{\leftmark}}{}{\footnotesize\thepage}
\makeheadrule{meuestilo}{\textwidth}{0.3pt}

% Redefine a marca para que \leftmark contenha o rótulo completo (e.g., "1. TÍTULO", "APÊNDICE A")
\renewcommand{\chaptermark}[1]{\markboth{\thechapter.\quad #1}{}}
\renewcommand{\sectionmark}[1]{\markright{\thechapter.\quad #1}}
% Customiza a página de abertura do capítulo para mostrar APENAS o número.
\aliaspagestyle{chapter}{plain}
%===========================================

% ===========================================
% IMPLEMENTAÇÃO IA: MACROS DE APÊNDICE/ANEXO
% Garante o reset do contador e a marcação correta do cabeçalho.
% ===========================================
\newcommand{\iniciaApendices}{
    \chapterstyle{estiloNaoNumerado}
    \cleardoublepage
    \setcounter{chapter}{0}
    \renewcommand{\thechapter}{\Alph{chapter}}
    \renewcommand{\chaptername}{APÊNDICE}
    \chapterstyle{CapitulosNumerado}
    % Redefine a marcação do cabeçalho para incluir 'APÊNDICE'
    \renewcommand{\chaptermark}[1]{\markboth{\chaptername\ \thechapter.\quad #1}{}}
    \renewcommand{\sectionmark}[1]{\markright{\chaptername\ \thechapter.\quad #1}}
}

\newcommand{\iniciaAnexos}{
    \chapterstyle{estiloNaoNumerado}
    \cleardoublepage
    \setcounter{chapter}{0}
    \renewcommand{\thechapter}{\Alph{chapter}}
    \renewcommand{\chaptername}{ANEXO}
    \chapterstyle{CapitulosNumerado}
    % Redefine a marcação do cabeçalho para incluir 'ANEXO'
    \renewcommand{\chaptermark}[1]{\markboth{\chaptername\ \thechapter.\quad #1}{}}
    \renewcommand{\sectionmark}[1]{\markright{\chaptername\ \thechapter.\quad #1}}
}
%===========================================


\begin{document}

% ===========================================
% INSTRUÇÃO IA: Configuração Inicial de Numeração e Estilo
% frontmatter: Numeração invisível, estilo de capítulo centralizado para pré-textuais.
% ===========================================
\frontmatter
\pagestyle{empty} % Garante que não há numeração visível em folha e contracapa
\chapterstyle{estiloNaoNumerado} % Define o estilo simplificado para pré-textuais

\folha
\contra

% Inicia a numeração das páginas pré-textuais, mas continua 'invisível' (sem printar)
% Por padrão memoir usa \pagenumbering{roman} em \frontmatter. Vamos garantir árabe,
% mas só se tornará visível em \mainmatter.
\pagenumbering{arabic}
\setcounter{page}{1}

%===========================================
% INSTRUÇÃO IA: ELEMENTOS PRÉ-TEXTUAIS
% Uso do comando modular \capituloCentralizado.
%===========================================
\capituloCentralizado{Agradecimento}
\capituloCentralizado{Resumo}
\capituloCentralizado{Lista de Figuras}
\capituloCentralizado{Lista de Tabelas}
\capituloCentralizado{Lista de Abreviaturas e Siglas}
\capituloCentralizado{Lista de Símbolos}

%===========================================
% INSTRUÇÃO IA: CONFIGURAÇÃO DO SUMÁRIO (CLASSE MEMOIR)
%===========================================
\renewcommand{\cfttoctitlefont}{\hfill\MakeUppercase\bfseries\Huge} % Título em CAIXA ALTA e negrito
\renewcommand{\cftaftertoctitle}{\hfill\par\vspace{2cm}} % Espaçamento horizontal (centralizado)
\setlength{\cftchapnumwidth}{2cm} % Espaçamento para o número do capítulo
\tableofcontents* % Supressão do título "SUMÁRIO" no índice. O título foi inserido manualmente.

% ===========================================
% Título "SUMÁRIO" centralizado manualmente
% O \tableofcontents* não gera o título, mas a classe memoir
% já cuida da formatação das entradas do Sumário.
% ===========================================
\chapterstyle{estiloNaoNumerado}
\chapter*{SUMÁRIO}
\chapterstyle{CapitulosNumerado} % Retorna para o estilo numerado (apesar de ainda estarmos em \frontmatter)


%===========================================
% INSTRUÇÃO IA: TRANSIÇÃO PARA A PARTE TEXTUAL (CRÍTICO)
% Numeração ARÁBICA (visível) e contínua.
% Ativa o estilo de cabeçalho 'meuestilo' e o estilo de capítulo 'CapitulosNumerado'.
% ===========================================
\mainmatter
\pagestyle{meuestilo}
\chapterstyle{CapitulosNumerado}
% Garante que a primeira numeração seja '1' (ou a próxima página, se for após o Sumário)

% A Introdução deve ser \chapter* (sem número) + \addcontentsline (para Sumário).
% Usamos \capituloCentralizado para garantir o estilo correto e entrada no TOC.
\capituloCentralizado{Introdução}

% Contagem de Capítulos DEVE continuar sequencialmente (1, 2, 3...)
\part{Preparação do Relatório}
\chapter{Desenvolvimento}
\section{Seção 1.1}
\section{Seção 1.2}
\part{Resultado do Relatório}
\chapter{Conclusão}
\section{Seção 2.1}

%===========================================
% INSTRUÇÃO IA: ELEMENTOS PÓS-TEXTUAIS (CRÍTICO)
%
% 1. Transição Pós-Textual: \backmatter
% 2. Referências: \chapter* (usando \capituloCentralizado).
% 3. Apêndice/Anexo: comandos customizados \iniciaApendices / \iniciaAnexos.
%===========================================
\backmatter
% Garante que as páginas de referência e outras não tenham número ou cabeçalho se o estilo de capítulo não for usado.

\capituloCentralizado{Referência}

% -------------------------------------------
% APÊNDICES
% -------------------------------------------
\iniciaApendices % Define \thechapter como \Alph, reseta o contador e define \chaptername
\capituloCentralizado{APÊNDICES} % Título divisório (sem numeração de página no TOC)
\chapter{Título do Apêndice A} % Deve gerar APÊNDICE A
\section{Seção A.1}
\chapter{Título do Apêndice B} % Deve gerar APÊNDICE B
\section{Seção B.1}

% -------------------------------------------
% ANEXOS
% -------------------------------------------
\iniciaAnexos % Define \thechapter como \Alph, reseta o contador e define \chaptername
\capituloCentralizado{ANEXOS} % Título divisório (sem numeração de página no TOC)
\chapter{Título do Anexo C} % Deve gerar ANEXO A (contador recomeçou)
\section{Seção A.1}
\chapter{Título do Anexo D} % Deve gerar ANEXO B
\section{Seção B.1}

\end{document}